\documentclass[12pt,a4paper]{article}
\usepackage[utf8]{inputenc}
\usepackage[T1]{fontenc}
\usepackage[ngerman]{babel}
\usepackage{amsmath,amssymb,amsthm}
\usepackage{geometry}
\usepackage{hyperref}
\usepackage{enumitem}
\usepackage{mathrsfs}
%\usepackage{fraktur}  % Removed - fraktur support provided by amssymb, mathfrak available
\usepackage{array}
\usepackage{booktabs}
\geometry{margin=2.5cm}

\theoremstyle{plain}
\newtheorem{satz}{Satz}[section]
\newtheorem{theorem}{Theorem}[section]
\newtheorem{lemma}{Hilfssatz}[section]
\newtheorem{hilfssatz}{Hilfssatz}[section]
\newtheorem{korollar}{Korollar}[section]
\newtheorem{folgerung}{Folgerung}[section]

\theoremstyle{definition}
\newtheorem{definition}{Definition}[section]
\newtheorem{bezeichnung}{Bezeichnung}[section]

\theoremstyle{remark}
\newtheorem{bemerkung}{Bemerkung}[section]
\newtheorem{beispiel}{Beispiel}[section]

\title{Lehrbuch der Algebra --- Erster Band, Teil 3}
\author{Heinrich Weber}
\date{}

\begin{document}
\maketitle

\section*{Dritter Abschnitt. Fortsetzung. Kreistheilungskörper.}

\subsection*{\S. 21. Kreistheilungskörper von gegebener Gruppe. (Fortsetzung)}

\noindent
Wenn die Primzahl $q$ nur einfach in $m$ aufgeht und $k$ durch
$Q$ theilbar ist, so ist zwar $k$ nicht primär; es gehört aber trotzdem
die Periode $\eta$ zu der Gruppe $k$.

Es kann aber in diesem Falle $\eta$ durch Einheitswurzeln von
niedrigerem Grade dargestellt werden. In folgender Weise lässt
sich diese Darstellung finden. Alle Zahlen von $Q$ sind von der
Form $1 + hm'$, worin $h$ die Reihe der Zahlen $0,1,\ldots q-1$
durchläuft, mit Ausnahme des einzigen Werthes, für den
\begin{equation}\tag{18}
1 + hm' \equiv kq
\end{equation}
durch $q$ theilbar wird. Setzt man also unter der Voraussetzung,
dass $k$ durch $Q$ theilbar ist,
\[
\eta = Q_1 + Q_2 + Q_3 + \cdots,
\]
so folgt, dass jede Zahl $a$ in $k\eta$ von der Form ist:
\begin{equation}\tag{19}
a \equiv a'(1 - hm') \pmod{m},
\end{equation}
worin $a'$ die Reihe der Zahlen $1, a_1, a_2, \ldots$ durchläuft. Daraus
ist noch zu schliessen, dass die Reste der Zahlen $1, a', a''\ldots$
nach dem Modul $m'$ eine Gruppe bilden. Nimmt man nun die
Summe über alle Werthe $h = 0, 1,\ldots q-1$ und nimmt dann
den einen durch (18) bestimmten Werth wieder weg, so folgt:
\[
\eta = \sum_{h} r^{a'} - \sum_{h} r^{a'kq},
\]
oder, da $\sum r^{a'h\cdot m'} = 0$ ist:
\[
\eta = -\sum_{h} r^{a'kq}.
\]
Nun ist $r^q$ eine $m'^{\text{te}}$ Einheitswurzel, und $\eta$ also, vom Vorzeichen abgesehen, gleich einer aus $m'^{\text{ten}}$ Einheitswurzeln gebildeten Periode.\footnote{Hier ist die Abhandlung von Fuchs zu vergleichen: Ueber die Perioden, welche aus den Wurzeln der Gleichung $w^n = 1$ gebildet sind, wenn $n$ eine zusammengesetzte Zahl ist. Crelle's Journal, Bd.~61 (1863).}

\subsection*{\S. 21. Kreistheilungskörper mit gegebener Gruppe}

Im Vorhergehenden hat sich also ergeben, dass jede Kreistheilungsperiode
\[
\eta = \sum_{\alpha} r^{a},
\]
in der $r$ eine primitive $m^{\text{te}}$ Einheitswurzel ist und $\alpha$ die Zahlen
einer Gruppe $\mathfrak{A}$ durchläuft, wenn $k$ ein primärer Theiler von $M$
ist, die Wurzel einer Abel'schen Gleichung ist, deren Gruppe mit
der zu $\mathfrak{A}$ reciproken Gruppe $\mathfrak{B}$ isomorph ist.

Die Aufgabe, die wir jetzt stellen und lösen wollen, ist folgende:

\medskip\noindent\textbf{I.} Es soll der Modul $m$ und die Gruppe $\mathfrak{A}$ auf alle mögliche Arten so bestimmt werden, dass $\eta$ ein beliebig gegebenes System von Invarianten hat.

Oder was dasselbe ist:

Es sollen alle Kreistheilungskörper von gegebener Gruppe bestimmt werden.

\medskip
Wir machen dabei immer die Voraussetzung, dass $k$ ein primärer Theiler von $\mathfrak{A}$ sein soll, weil wir sonst einen und denselben Körper mehrmals erhalten würden. Die Aufgabe zerfällt in zwei Theile, nämlich:

\medskip\noindent\textbf{II.} Welche Moduln $m$ sind geeignet, eine Gruppe $\mathfrak{B}$ von gegebenen Invarianten zu erzeugen?

\medskip\noindent\textbf{III.} Wie findet man diese Gruppe $\mathfrak{B}$ und die zugehörige Gruppe $\mathfrak{A}$?

\medskip
Wir müssen bei dieser Untersuchung die Bezeichnung gegen das Frühere etwas ändern, damit die Uebersichtlichkeit nicht verloren geht. Die exceptionelle Stellung der Primzahl $2$, die bei den bisherigen Betrachtungen immer berücksichtigt werden musste, machte eine etwas umständliche Bezeichnung nothwendig. Diese Unterscheidung ist im Folgenden nicht mehr in dem Maasse notwendig, und darum lässt sich die Bezeichnung jetzt vereinfachen.

Wir bezeichnen die Indices einer Zahl $a$ aus der Gruppe $\mathfrak{A}$ mit
\[
\alpha_1, \alpha_2, \ldots \alpha_u \quad \text{(Indices von $a$)}
\]
und die entsprechenden Indexmoduln mit
\[
c_1, c_2, \ldots c_u \quad \text{(Indexmoduln)},
\]
so dass $u$ gleich der Anzahl der in $m$ aufgehenden Primzahlen,
oder wenn $m$ durch $8$ theilbar ist, um $1$ grösser ist.

Es seien ferner
\[
\omega_1, \omega_2, \ldots \omega_u
\]
primitive Einheitswurzeln der Grade $c_1, c_2, \ldots c_u$. Sind nun die Indices einer Zahl $b$ in der zu $\mathfrak{A}$ reciproken Gruppe $\mathfrak{B}$
\[
\beta_1, \beta_2, \ldots \beta_u \quad \text{(Indices von $b$)},
\]
so ist die Beziehung zwischen den $\alpha$ und den $\beta$ durch die Gleichung
\begin{equation}\tag{1}
\omega_1^{\alpha_1\beta_1} \omega_2^{\alpha_2\beta_2} \cdots \omega_u^{\alpha_u\beta_u} = 1
\end{equation}
ausgedrückt.

Die Invarianten der Gruppe $\mathfrak{B}$, die nach \S.~10 lauter Primzahlpotenzen sind, wollen wir mit
\[
i_1, i_2, \ldots i_v \quad \text{(Invarianten von $\mathfrak{B}$)}
\]
bezeichnen. Wir nehmen diese Invarianten als beliebig gegebene Primzahlpotenzen an und bezeichnen mit $J$ ihr kleinstes gemeinschaftliches Vielfache, so dass
\begin{equation}\tag{2}
J = i_1 i_2 \cdots i_v
\end{equation}
gesetzt werden kann.

Ausserdem soll $k$ als primärer Theiler von $J$ vorausgesetzt werden, was nach \S.~19, 8. mit der Bedingung gleichbedeutend ist:

\medskip\noindent\textbf{IV.} Keiner der Indices $\beta$ soll für alle Zahlen $b$, wenn die entsprechende Primzahl $q$ mehrmals in $m$ aufgeht, durch $q$, oder wenn $q$ nur einmal in $m$ aufgeht, durch $q-1$ theilbar sein.

\medskip
Nach der Definition der Invarianten muss die Gruppe $\mathfrak{B}$ eine Basis haben, deren Elemente
\[
g_1, g_2, \ldots g_v \quad \text{(Basis von $\mathfrak{B}$)}
\]
von den Graden $i_1, i_2, \ldots i_v$ sind, so dass jede Zahl $b$ einmal und nur einmal in der Form
\begin{equation}\tag{3}
b \equiv g_1^{z_1} g_2^{z_2} \cdots g_v^{z_v} \pmod{m}
\end{equation}
enthalten ist, wenn die Exponenten $z_1, z_2, \ldots z_v$ je ein volles Restsystem nach den Moduln $i_1, i_2, \ldots i_v$ durchlaufen.

Alle Grössen von der Form (3) bilden, wenn $g_1, g_2, \ldots g_v$ beliebige Zahlen sind, bei unbeschränkter Veränderlichkeit der ganzzahligen Exponenten $z$ gewiss eine Gruppe. Sollen aber die $g_i$ die Elemente einer Basis dieser Gruppe, und ihre Grade $i_j$ sein, so darf $1 \equiv g_1^{z_1} g_2^{z_2} \cdots g_v^{z_v} \pmod{m}$ nur dann erfüllt sein, wenn $z_j$ durch $i_j$ theilbar ist, also:

\medskip\noindent\textbf{1.} Die nothwendige und hinreichende Bedingung dafür, dass $g_1, g_2, \ldots g_v$ die Elemente einer Basis einer Gruppe $\mathfrak{B}$ von den Graden $i_1, i_2, \ldots i_v$ seien, ist die, dass die Congruenz
\begin{equation}\tag{4}
g_1^{z_1} g_2^{z_2} \cdots g_v^{z_v} \equiv 1 \pmod{m}
\end{equation}
dann und nur dann besteht, wenn zugleich die Congruenzen
\begin{equation}\tag{5}
z_1 \equiv 0 \pmod{i_1}, \quad z_2 \equiv 0 \pmod{i_2}, \quad \ldots \quad z_v \equiv 0 \pmod{i_v}
\end{equation}
erfüllt sind.

\medskip
Hiernach ist also $J$ die kleinste positive Zahl, die für alle $b$ der Congruenz
\[
b^J \equiv 1 \pmod{m}
\]
genügt, oder die kleinste positive Zahl, für die für alle Systeme der $\beta$
\begin{equation}\tag{6}
J\beta_1 \equiv 0 \pmod{c_1}, \quad J\beta_2 \equiv 0 \pmod{c_2}, \quad \ldots \quad J\beta_u \equiv 0 \pmod{c_u}.
\end{equation}

Daraus ergeben sich die ersten Schlüsse über die Zusammensetzung der Zahl $m$.

\medskip\noindent\textbf{a)} Eine ungerade Primzahl $q$, die nicht in $J$ enthalten ist, kann nur einfach in $m$ aufgehen.

Denn ist $m$ durch $q^x$ theilbar, so ist eine der Grössen $c$, etwa $c_1 = \varphi(q^x) = q^{x-1}(q-1)$, also wenn $x > 1$ ist, so ist $c_1$ noch durch $q$ theilbar, und aus (6) folgt dann, dass alle $\beta_i$ durch $q$ theilbar sein müssten, was der Voraussetzung IV. widerspricht.

\medskip\noindent\textbf{b)} Eine ungerade Primzahl $q$, die in $J$ aufgeht, kann höchstens einmal mehr in $m$, als in $J$ enthalten sein.

Denn man kann ebenschliessen, dass, wenn $c_1 = q^{x-1}(q-1)$, und $J$ nicht durch $q^{x-1}$ theilbar ist, alle $\beta_i$ durch $q$ theilbar sein müssten.

\medskip\noindent\textbf{c)} Ist $J$ ungerade, so muss auch $m$ ungerade sein.

Denn ist $m$ durch eine Potenz von $2$, also mindestens durch $4$ theilbar, so ist der der Zahl $2$ entsprechende Indexmodul gerade, und die entsprechenden $\beta$ müssten nach (6) alle gerade sein, entgegen der Forderung IV.

\medskip\noindent\textbf{d)} Ist $J$ durch eine Potenz von $2$ theilbar, so kann $m$ den Factor $2$ höchstens zweimal öfter enthalten, als $J$.

Denn ist $m$ durch $2^x$ theilbar und $x > 2$, so sind zwei der Indexmoduln $c$ gleich $2^{x-2}$ oder $2^{x-1}$. Wäre also $J$ nicht durch $2^{x-2}$ theilbar, so müssten alle $\beta_i$ durch $2$ theilbar sein, was wieder gegen die Forderung IV ist.

\medskip\noindent\textbf{e)} Wenn $q$ eine einfach in $m$ aufgehende Primzahl ist, so muss $q-1$ wenigstens durch eine der in $J$ aufgehenden Primzahlen theilbar sein.

Denn wäre $c_1 = q-1$ relativ prim zu $J$, so müssten alle $\beta_i$ durch $q-1$ theilbar sein, im Widerspruch mit IV. Die weiteren Bedingungen für $m$ ergeben sich später.

\subsection*{Bestimmung der Gruppe $\mathfrak{B}$}

Die Gruppe $\mathfrak{B}$ ist bestimmt, wenn ihre Basis $g_1, g_2, \ldots g_v$ gegeben ist. Die Elemente der Basis wollen wir durch ihre Indices bestimmen, und führen also folgende Bezeichnung ein. Es seien
\begin{equation}\tag{7}
\begin{array}{cccc}
\gamma_{1,1}, & \gamma_{1,2}, & \ldots & \gamma_{1,u} \quad \text{die Indices von $g_1$}, \\
\gamma_{2,1}, & \gamma_{2,2}, & \ldots & \gamma_{2,u} \quad \text{die Indices von $g_2$}, \\
\vdots & \vdots & & \vdots \\
\gamma_{v,1}, & \gamma_{v,2}, & \ldots & \gamma_{v,u} \quad \text{die Indices von $g_v$}. \\
\end{array}
\end{equation}

Die Indices $\beta_1, \beta_2, \ldots \beta_u$ eines beliebigen Elementes $b$ erhält man dann nach (3) in der Form:
\begin{equation}\tag{8}
\begin{aligned}
\beta_1 &\equiv \gamma_{1,1} z_1 + \gamma_{2,1} z_2 + \cdots + \gamma_{v,1} z_v \pmod{c_1}, \\
\beta_2 &\equiv \gamma_{1,2} z_1 + \gamma_{2,2} z_2 + \cdots + \gamma_{v,2} z_v \pmod{c_2}, \\
&\vdots \\
\beta_u &\equiv \gamma_{1,u} z_1 + \gamma_{2,u} z_2 + \cdots + \gamma_{v,u} z_v \pmod{c_u},
\end{aligned}
\end{equation}
und wir fragen nun nach den nothwendigen und hinreichenden Bedingungen für die Zahlen $\gamma_{j,i}$, damit $g_1, g_2, \ldots g_v$ die Basis einer Gruppe $\mathfrak{B}$ von den verlangten Eigenschaften ist. Nach 1. ist hierfür zuerst nothwendig:

\medskip\noindent\textbf{2.} Die Zahlen $\gamma_{j,i}$ müssen so beschaffen sein, dass die Congruenzen
\begin{equation}\tag{9}
\begin{aligned}
\gamma_{1,1} z_1 + \gamma_{2,1} z_2 + \cdots + \gamma_{v,1} z_v &\equiv 0 \pmod{c_1}, \\
\gamma_{1,2} z_1 + \gamma_{2,2} z_2 + \cdots + \gamma_{v,2} z_v &\equiv 0 \pmod{c_2}, \\
&\vdots \\
\gamma_{1,u} z_1 + \gamma_{2,u} z_2 + \cdots + \gamma_{v,u} z_v &\equiv 0 \pmod{c_u}
\end{aligned}
\end{equation}
dann und nur dann erfüllt sind, wenn zugleich die Congruenzen
\begin{equation}\tag{10}
z_1 \equiv 0 \pmod{i_1}, \quad z_2 \equiv 0 \pmod{i_2}, \quad \ldots \quad z_v \equiv 0 \pmod{i_v}
\end{equation}
bestehen.

\medskip
Ist diese Bedingung erfüllt, so ist gewiss $g_1, g_2, \ldots g_v$ die Basis einer Gruppe mit den Invarianten $i_1, i_2, \ldots i_v$. Wenn aber auch noch die in IV. aufgestellte Forderung erfüllt sein soll, dann folgt noch weiter:

\medskip\noindent\textbf{3.} Ist $q_h$ eine in $m$ aufgehende Primzahl, der der Index $\beta_h$ entspricht, so dürfen die $v$ Grössen
\[
\gamma_{1,h}, \gamma_{2,h}, \ldots \gamma_{v,h}
\]
wenn $q_h$ mehrfach in $m$ aufgeht, nicht alle durch $q_h$, und wenn $q_h$ nur einmal in $m$ aufgeht, nicht alle durch $q_h-1$ theilbar sein.

\medskip
Die Bedingungen 2., 3. sind dann also die nothwendigen und hinreichenden, und es kommt jetzt nur noch darauf an, sie in eine Form zu bringen, dass die Möglichkeit ihrer Erfüllung beurtheilt werden kann.

Wir bezeichnen mit $\delta_{k,h}$ den grössten gemeinschaftlichen Theiler von $i_k$ und $c_h$, und setzen
\[
i_k = \delta_{k,h} \cdot i_{k,h}, \qquad c_h = \delta_{k,h} \cdot c_{k,h},
\]
\[
k = 1, 2, \ldots v; \quad h = 1, 2, \ldots u,
\]
so dass $i_{k,h}$ und $c_{k,h}$ relativ prim sind. Es ist dann zu bemerken, dass die Grössen
\begin{equation}\tag{11}
\delta_{k,h}, \quad i_{k,h}, \quad c_{k,h}
\end{equation}
durch $m$ und die Invarianten $i$ völlig bestimmt sind. Nach der Definition der Invarianten müssen die $\delta_{k,h}$ und $i_{k,h}$ Primzahlpotenzen sein.

Aus den Voraussetzungen, die wir in a) bis d) über die Zahl $m$ gemacht haben, ergiebt sich eine Folgerung für die $c_h$, die wir hervorheben müssen.

\medskip\noindent\textbf{4.} Ist $q^x$ einer der Factoren von $m$, der dem Index $\beta_h$ entspricht, ist also $c_h = q^{x-1}(q-1)$, oder wenn $q=2$ und $x > 2$ ist, $c_h = 2^{x-2}$, und wenn $x=2$ ist, $c_h = 2$, so können, wenn $x > 1$ ist, die Zahlen
\begin{equation}\tag{12}
c_{1,h}, c_{2,h}, \ldots c_{v,h}
\end{equation}
nicht alle durch $q$, und wenn $x=1$ ist, nicht alle durch $q-1$ theilbar sein.

\medskip
Wenn nämlich die Grössen (12) alle durch $q$ theilbar sind, so sind wegen $c_h = \delta_{k,h} c_{k,h}$ alle $\delta_{k,h}$ Potenzen von $q$. Da nun nach (11) $i_k = \delta_{k,h} i_{k,h}$ ist, so können die $i_k$ keine andere Primzahl als $q$ enthalten, und da die $i_k$ die Invarianten der Gruppe $\mathfrak{B}$ sind, so können auch diese selbst nur durch $q$ theilbar sein. Ist also $x > 1$, so enthält $m$ die Primzahl $q$, und es folgt aus (6), dass alle $\beta_h$ durch $q$ theilbar sein müssten, im Widerspruch mit IV.



\medskip
Wir wollen nun annehmen, dass die Bedingung 4. erfüllt sei, und wollen zeigen, dass dann die Bedingungen 2. und 3. immer erfüllt werden können.

Zu dem Ende setzen wir
\begin{equation}\tag{13}
\gamma_{k,h} \equiv \varepsilon_{k,h} c_{k,h} \pmod{c_h},
\end{equation}
wobei die $\varepsilon_{k,h}$ noch näher zu bestimmende ganze Zahlen sind, und es ist dann
\[
\frac{\gamma_{k,h}}{c_h} = \frac{\varepsilon_{k,h}}{\delta_{k,h}}.
\]
Die Congruenzen (9) verwandeln sich dadurch in
\[
\frac{z_1 \varepsilon_{1,h}}{\delta_{1,h}} + \frac{z_2 \varepsilon_{2,h}}{\delta_{2,h}} + \cdots + \frac{z_v \varepsilon_{v,h}}{\delta_{v,h}} \equiv 0 \pmod{1},
\]
worin $h = 1, 2, \ldots u$ ist.

Bringt man die Glieder auf den gemeinschaftlichen Nenner
\[
\Delta_h = [\delta_{1,h}, \delta_{2,h}, \ldots \delta_{v,h}],
\]
so erhält man:
\begin{equation}\tag{14}
\frac{z_1 \varepsilon_{1,h} \Delta_h}{\delta_{1,h}} + \frac{z_2 \varepsilon_{2,h} \Delta_h}{\delta_{2,h}} + \cdots + \frac{z_v \varepsilon_{v,h} \Delta_h}{\delta_{v,h}} \equiv 0 \pmod{\Delta_h}.
\end{equation}

Die Coefficienten von $z_1, z_2, \ldots z_v$ sind durch $\Delta_h/\delta_{k,h}$ theilbar. Nun ist $\Delta_h/\delta_{k,h}$ relativ prim zu $i_{k,h}$, aber ein Theiler von $\delta_{k,h}$, also auch von $i_k$. Daher kann man, wenn $\varepsilon_{k,h}$ durch $i_{k,h}$ theilbar ist, setzen:
\[
\varepsilon_{k,h} = \frac{\delta_{k,h}}{\Delta_h} \cdot \zeta_{k,h},
\]
worin $\zeta_{k,h}$ eine beliebige ganze Zahl ist.

Dann geht (14) über in
\begin{equation}\tag{15}
\zeta_{1,h} z_1 + \zeta_{2,h} z_2 + \cdots + \zeta_{v,h} z_v \equiv 0 \pmod{\Delta_h}.
\end{equation}

Nun ist $\Delta_h$ ein Theiler von $i_1 i_2 \cdots i_v$. Wir wollen nun die Zahlen $\zeta_{k,h}$ so wählen, dass die Congruenzen (15) dann und nur dann bestehen, wenn $z_k \equiv 0 \pmod{i_k}$ ist. Dies wird erreicht, wenn man die Determinante
\begin{equation}\tag{16}
\begin{vmatrix}
\zeta_{1,1} & \zeta_{2,1} & \cdots & \zeta_{v,1} \\
\zeta_{1,2} & \zeta_{2,2} & \cdots & \zeta_{v,2} \\
\vdots & \vdots & & \vdots \\
\zeta_{1,v} & \zeta_{2,v} & \cdots & \zeta_{v,v}
\end{vmatrix}
\end{equation}
gleich $\pm 1$ setzt.

Die Zahlen $\zeta_{k,h}$ können wir z.B. durch die Matrix der Einheitscoefficienten bestimmen:
\[
\zeta_{k,h} = \begin{cases} 1 & \text{wenn } k = h, \\ 0 & \text{wenn } k \neq h, \end{cases} \qquad (k,h = 1, 2, \ldots v)
\]
und dann ist
\[
\varepsilon_{k,h} = \frac{\delta_{k,h}}{\Delta_h} \zeta_{k,h},
\]
wobei $\varepsilon_{k,h} = 0$ ist für $k \neq h$, und $\varepsilon_{k,k} = \delta_{k,h}/\Delta_h$.

Nun muss noch die Bedingung 3. erfüllt werden. Die $v$ Grössen
\[
\gamma_{1,h}, \gamma_{2,h}, \ldots \gamma_{v,h}
\]
dürfen, wenn $q_h$ mehrfach in $m$ aufgeht, nicht alle durch $q_h$, und wenn $q_h$ nur einmal in $m$ aufgeht, nicht alle durch $q_h-1$ theilbar sein.

Nun ist nach (13)
\[
\gamma_{k,h} \equiv \varepsilon_{k,h} c_{k,h} = \zeta_{k,h} \frac{\delta_{k,h} c_{k,h}}{\Delta_h} = \zeta_{k,h} \frac{c_h}{\Delta_h} \pmod{c_h}.
\]
Da $\zeta_{k,h} = 0$ für $k \neq h$ und $\zeta_{h,h} = 1$ ist, so reduciren sich die Grössen $\gamma_{k,h}$ auf die eine Zahl $c_h/\Delta_h$, und die Bedingung 3. ist erfüllt, wenn diese Zahl nicht durch $q_h$ (bez. $q_h-1$) theilbar ist.

Es ist aber $\Delta_h$ das kleinste gemeinschaftliche Vielfache der zu einer und derselben Primzahl $q$ gehörigen $\delta_{k,h}$. Wenn also $x > 1$ ist, so ist $\Delta_h = q^\lambda$, wo $\lambda$ die grösste in den $\delta_{k,h}$ vorkommende Potenz ist. Dann ist $c_h/\Delta_h = q^{x-1-\lambda}(q-1)$, und dies ist nur dann durch $q$ theilbar, wenn $x-1-\lambda > 0$, d.h. wenn $\Delta_h < q^{x-1}$. Wenn aber $x=1$ ist, so ist $c_h = q-1$, und $\Delta_h$ theilt $q-1$.

\medskip\noindent\textbf{5.} Die nothwendige und hinreichende Bedingung dafür, dass die Bedingungen 2. und 3. erfüllt werden können, ist also die, dass für jedes $h$, dem eine mehrfache Primzahl $q_h$ entspricht, die zugehörige Zahl $c_h/\Delta_h$ nicht durch $q_h$, und für jedes $h$, dem eine einfache Primzahl entspricht, $c_h/\Delta_h$ nicht durch $q_h-1$ theilbar sei.

\medskip
Dies ist aber genau die Bedingung, die in 4. ausgesprochen ist. Denn wenn alle $c_{k,h}$ durch $q$ theilbar sind, so sind alle $\delta_{k,h}$ durch $q$ theilbar, und $\Delta_h$ ist durch $q$ theilbar, also $c_h/\Delta_h$ nicht durch $q^{x-1}(q-1)/q^{x-1} = q-1$, sondern nur dann durch $q$, wenn $\Delta_h < q^{x-1}$.

\medskip\noindent\textbf{6.} Wir können also die $\gamma_{k,h}$ in der einfachsten Weise durch die Formeln
\begin{equation}\tag{17}
\gamma_{k,h} \equiv \frac{c_h}{\Delta_h} \zeta_{k,h} \pmod{c_h}
\end{equation}
bestimmen, worin die $\zeta_{k,h}$ ganze Zahlen mit der Determinante $\pm 1$ sind.

\medskip
Diese Bestimmung genügt aber nicht allen Forderungen. Wir müssen nämlich noch die Bedingung erfüllen, dass die Gruppe $\mathfrak{A}$, die zu der so bestimmten Gruppe $\mathfrak{B}$ reciprok ist, einen primären Theiler von $J$ habe. Diese Bedingung lässt sich nun folgendermaßen aussprechen:

\medskip\noindent\textbf{7.} Es muss die Anzahl $g$ der durch eine Primzahl $p$ theilbaren Invarianten gleich oder kleiner als die Anzahl $w$ der Indexmoduln sein, und es muss sich aus der Matrix
\begin{equation}\tag{24}
\begin{pmatrix}
c_{1,1} & c_{2,1} & c_{3,1} & \cdots & c_{g,1} \\
c_{1,2} & c_{2,2} & c_{3,2} & \cdots & c_{g,2} \\
\vdots & \vdots & \vdots & & \vdots \\
c_{1,w} & c_{2,w} & c_{3,w} & \cdots & c_{g,w}
\end{pmatrix}
\end{equation}
wenigstens eine durch $p$ untheilbare Determinante von $g$ Reihen bilden lassen.

\medskip
Dies involvirt zunächst eine Forderung für die $c_h$, d.h. also in letzter Instanz eine Bedingung für die Zahl $m$.

Alle Glieder irgend einer $g$-reihigen Determinante der Matrix (24) enthalten einen Factor der Form
\[
c_{h_1,1} c_{h_2,2} \cdots c_{h_g,g},
\]
worin $h_1, h_2, \ldots h_g$ irgend eine Combination von $g$ verschiedenen Zahlen der Reihe $1, 2, \ldots w$ bedeuten. Wären nun alle diese Producte durch $p$ theilbar, so wären sicher alle Determinanten der Matrix (24) von $g$ Reihen auch durch $p$ theilbar, und es muss also wenigstens eine solche Combination geben, die nicht durch $p$ theilbar ist. Weil aber die $i_k$ nur Potenzen von $p$ sein können, so muss
\[
c_{h_1,1} = 1, \quad c_{h_2,2} = 1, \quad \ldots \quad c_{h_g,g} = 1
\]
sein. Geht man nun auf die Bedeutung der $c_{k,h}$ in (11) zurück, so können wir die letzte Forderung für den Grad $m$ folgendermaassen ausdrücken:

\medskip\noindent\textbf{f)} Wenn eine Primzahl $p$ in $g$ Invarianten $i_1, i_2, \ldots i_g$ aufgeht, so muss die Anzahl $w$ der Indexmoduln gleich oder grösser als $g$ sein, und es müssen $g$ dieser Indexmoduln $c_{h_1}, c_{h_2}, \ldots c_{h_g}$ so auswählen lassen, dass $c_{h_1}$ durch $i_1$, $c_{h_2}$ durch $i_2$, $\ldots$, $c_{h_g}$ durch $i_g$ theilbar ist.\footnote{Der Nachweis der Thatsache, dass dieser Forderung immer auf unendlich viele Arten entsprochen werden kann, stützt sich auf den Satz der Zahlentheorie, dass unter den Zahlen einer arithmetischen Progression, deren Anfangsglied und Differenz ohne gemeinsamen Theiler sind, unendlich viele Primzahlen vorkommen, ein Satz, der bis jetzt nur von Dirichlet mit Anwendung der Integralrechnung bewiesen ist. (Abhandlungen der Berliner Akademie 1837; Dirichlet's Werke, Bd.~I, Nr.~21; Dirichlet-Dedekind, Vorlesungen über Zahlentheorie, Supplement~VI; Bachmann, Analytische Zahlentheorie, Abschn.~IV). Behalten wir die oben benutzte Bezeichnung bei, so kann man, um für ein gegebenes $p$ der Forderung f) zu genügen, $p^{x+1}$ als Factor in $m$ aufnehmen. Dann aber müssen, wenn $g > 1$ ist, noch $g-1$ Primfactoren $q$ in $m$ vorkommen, die an Congruenzen von der Form $q \equiv 1 \pmod{i_k}\ldots$ gebunden sind.}

\medskip
Hiermit aber ist alles erschöpft, was wir von $m$ fordern müssen. Denn wenn diese Bedingung f) befriedigt ist, so brauchen wir z.~B. nur die $c_{1,h_1}, c_{2,h_2}, \ldots c_{g,h_g}$ durch $p$ untheilbar, die übrigen $c$ durch $p$ theilbar anzunehmen; dann ist die Determinante
\[
\begin{vmatrix}
c_{1,h_1} & c_{2,h_1} & \cdots & c_{g,h_1} \\
c_{1,h_2} & c_{2,h_2} & \cdots & c_{g,h_2} \\
\vdots & \vdots & & \vdots \\
c_{1,h_g} & c_{2,h_g} & \cdots & c_{g,h_g}
\end{vmatrix}
\equiv 1 \cdot c_{2,h_2} \cdots c_{g,h_g} \pmod{p},
\]
also sicher durch $p$ untheilbar.

Und bei dieser Annahme kann auch noch die Bedingung 6. befriedigt werden. Um dies einzusehen, bezeichne man mit $p^x$ irgend eine der Invarianten, etwa $i_1$, und mit $c_h$ den Indexmodul, der nach f) durch $p^x$ theilbar ist, so dass
\begin{equation}\tag{25}
c_h = p^x \cdot c_{h}'
\end{equation}
ist. Ist dann $q_h$ der zum Indexmodul $c_h$ gehörige Primtheiler von $m$, so ist $c_h$ entweder gleich $q_h$ oder ein Theiler von $q_h-1$.

Nach der zuletzt gemachten Annahme ist $c_{1,h}$ nicht durch $p$ theilbar, während $c_{2,h}, c_{3,h}, \ldots c_{g,h}$ durch $p$ theilbar sind. Es sollen dann die Producte
\[
c_{1,h} c_{h}', \quad c_{2,h} c_{h}', \quad \ldots \quad c_{g,h} c_{h}'
\]
nicht alle durch $q_h$, oder durch $q_h-1$ theilbar sein.

Ist zunächst $q_h$ ein einfacher Theiler von $m$, so ist $c_h = q_h-1$, und wegen (25) ist $c_{1,h}$ nicht durch $q_h-1$ theilbar. Man kann also $c_{1,h}$ noch so annehmen, dass $c_{1,h} c_{h}'$ nicht durch $q_h-1$ theilbar ist.

Ist aber $q_h$ ein $x$-facher Theiler von $m$ und $x > 1$, und ist $q_h = p$, so ist, wenn $p$ ungerade ist, $c_h = p^{x-1}(p-1)$, und nach der Voraussetzung b) z.~\S.~$x-1$. Wenn also $c_h$ durch $p^x$ theilbar sein soll, so muss $x = x-1$ sein, und (25) ergiebt $c_{h}' = p-1$. Ist aber $p = 2$, so ist $c_h = 2$ oder $= 2^{x-1}$, und es muss also nach der Voraussetzung d) $x = 1$ oder $= x-2$ sein, und es ist $c_{1,h} = 1$. Man kann also auch in diesen Fällen $c_{1,h}$ so annehmen, dass $c_{1,h} c_{h}'$ durch $q_h$ nicht theilbar wird.

Die Forderungen, die hiermit für die $c_h$ gestellt sind, bestehen also nur darin, dass sie durch gewisse Primzahlen theilbar, durch andere nicht theilbar sein sollen, und sind also noch auf unendlich viele Arten mit einander verträglich. Aber diese letzte Annahme über die $c_h$ hatte nur den Zweck, zu zeigen, dass die früheren allgemeineren Forderungen immer befriedigt werden können. Es kann noch viele andere Arten geben, ihnen zu genügen.

\subsection*{Uebersicht der Resultate}

Wenn es sich darum handelt, alle Kreistheilungskörper, und jeden nur einmal, und zwar auf möglichst einfache Art, durch Kreistheilungsperioden darzustellen, so verfahre man so:

Man nehme ein beliebiges System von Primzahlpotenzen als Invarianten
\[
i_1, i_2, \ldots i_v
\]
an, und wähle dann eine Zahl $m$ nach folgenden Bedingungen.

Wenn $p^x$ die höchste Potenz einer Primzahl $p$ ist, die unter den Invarianten vorkommt, so nehme man $p$ höchstens $(x+1)$-, oder, wenn $p = 2$ ist, $(x+2)$ mal in $m$ auf.

Eine Primzahl $q$, die gar nicht in den Invarianten aufgeht, darf nur einfach in $m$ aufgenommen werden; aber nur solche einfache Primfactoren $q$ darf $m$ haben, bei denen $q-1$ durch einen der Primfactoren $p$ der Invarianten theilbar ist. Also kann auch $p$ selbst nur dann einfach in $m$ vorkommen, wenn $p-1$ durch eines der anderen $p$ theilbar ist.

Ferner muss $m$ noch der Bedingung genügen, dass unter den Indexmoduln $c_1, c_2, \ldots c_u$ $g$ verschiedene, wie $c_{h_1}, c_{h_2}, \ldots c_{h_g}$ durch $i_1, i_2, \ldots i_g$ theilbar sind, wenn $i_1, i_2, \ldots i_g$ Potenzen von einer und derselben Primzahl sind. Dann bestimme man die Grössen $\delta_{k,h}$ und $c_{k,h}$ nach den Formeln (11) und theile die gegebenen Invarianten in Systeme ein, so dass alle Potenzen einer und derselben Primzahl in einem Systeme vereinigt sind.

Für jedes dieser Systeme bilde man die Matrix (24) und wähle die Zahlen $c_h$ so, dass aus jeder solchen Matrix eine Determinante von $g$ Reihen gebildet werden kann, die durch die betreffende Primzahl nicht theilbar ist, und dass nicht alle Producte
\[
c_{1,h} c_{h}', \quad c_{2,h} c_{h}', \quad \ldots \quad c_{g,h} c_{h}'
\]
durch $q_h$ oder durch $q_h-1$ theilbar werden. Dann setze man
\[
\gamma_{k,h} = c_h \cdot c_{k,h}
\]
und hat so nach (7) die Indices einer Basis einer Abel'schen Gruppe $\mathfrak{B}$, deren reciproke Gruppe $\mathfrak{A}$ ein primärer Theiler der ganzen Gruppe $\mathfrak{M}$ ist.



\subsection*{\S. 22. Bestimmung der Gruppe $\mathfrak{A}$}

Das letzte Ziel dieser Untersuchung ist nicht sowohl die Bestimmung der Gruppe $\mathfrak{B}$, als die der Gruppe $\mathfrak{A}$, aus der man direct die Kreistheilungsperioden $\eta = \sum r^a$ bilden kann.

Diese Aufgabe kann man auf folgende Art lösen: Nach (1) sind die Indices $\alpha$ der Zahlen in $\mathfrak{A}$ von den $\beta$ abhängig durch die Gleichung
\begin{equation}\tag{1}
\omega_1^{\alpha_1\beta_1} \omega_2^{\alpha_2\beta_2} \cdots \omega_u^{\alpha_u\beta_u} = 1,
\end{equation}
worin $\omega_1, \omega_2, \ldots \omega_u$ feste primitive Einheitswurzeln der Grade $c_1, c_2, \ldots c_u$ bedeuten. Versteht man aber unter $C$ das kleinste gemeinschaftliche Vielfache von $c_1, c_2, \ldots c_u$ und setzt
\begin{equation}\tag{2}
C = C_1 c_1 = C_2 c_2 = \cdots = C_u c_u,
\end{equation}
und bezeichnet mit $\Omega$ eine primitive $C^{\text{te}}$ Einheitswurzel, so kann man statt (1) auch setzen
\[
\Omega^{C_1 \alpha_1 \beta_1 + C_2 \alpha_2 \beta_2 + \cdots + C_u \alpha_u \beta_u} = 1,
\]
und bekommt daher für die $\alpha$ die Congruenz
\[
C_1 \alpha_1 \beta_1 + C_2 \alpha_2 \beta_2 + \cdots + C_u \alpha_u \beta_u \equiv 0 \pmod{C},
\]
die gleichbedeutend ist mit der Bedingung, dass die Summen
\[
\frac{C_1 \alpha_1 \beta_1}{c_1}, \quad \frac{C_2 \alpha_2 \beta_2}{c_2}, \quad \ldots \quad \frac{C_u \alpha_u \beta_u}{c_u}
\]
ganze Zahlen sein müssen. Wenn man darin für $\beta_i$ die Ausdrücke \S.~21, (8) substituirt, so folgt, dass auch
\[
\sum_h \frac{C_h \alpha_h \gamma_{k,h}}{c_h}
\]
ganze Zahlen sein müssen. Setzt man dann nach (15) und (11) \S.~21
\[
\varepsilon_{k,h} = c_{k,h} c_h = \delta_{k,h} c_{k,h} \cdot i_{k,h} = c_h \cdot i_{k,h},
\]
so folgt, dass
\[
\frac{C_h \alpha_h \gamma_{k,h}}{c_h} = C_h \alpha_h \varepsilon_{k,h}
\]
ganze Zahlen sein müssen, und diese Forderung ist gleichbedeutend mit dem System der Congruenzen
\[
C_h \alpha_h \varepsilon_{k,h} \equiv 0 \pmod{i_k},
\]
oder, ausführlicher geschrieben
\begin{equation}\tag{3}
\begin{aligned}
c_{1,1} \alpha_1 + c_{2,1} \alpha_2 + \cdots + c_{g,1} \alpha_u &\equiv 0 \pmod{i_1}, \\
c_{1,2} \alpha_1 + c_{2,2} \alpha_2 + \cdots + c_{g,2} \alpha_u &\equiv 0 \pmod{i_2}, \\
&\vdots \\
c_{1,v} \alpha_1 + c_{2,v} \alpha_2 + \cdots + c_{g,v} \alpha_u &\equiv 0 \pmod{i_v}.
\end{aligned}
\end{equation}

Durch diese Congruenzen sind, wie es sein muss, die Zahlen $\alpha_1, \alpha_2, \ldots \alpha_u$ nur nach den Moduln $c_1, c_2, \ldots c_u$ bestimmt, weil $c_h \cdot c_{k,h}$ durch $i_k$ theilbar ist, und um die Indices aller Zahlen der Gruppe $\mathfrak{A}$ zu finden, hat man einfach alle Lösungen der Congruenzen (3) aufzusuchen.

\newpage
\section*{Vierter Abschnitt.}
\section*{Cubische und biquadratische Abel'sche Körper.}

\subsection*{\S. 23. Cubische Kreistheilungskörper}

Die allgemeinen Untersuchungen, die in den letzten Paragraphen dargestellt sind, gestatten mannigfache Anwendungen auf specielle Fälle, von denen die einfachsten hier betrachtet werden sollen. Es ist dabei bemerkenswerth, dass schon die einfachsten Fälle, wenn man sie direct angreift, fast in dieselben Schwierigkeiten hineinführen, die wir durch die allgemeine Untersuchung überwunden haben.

Wir wollen zuerst die Aufgabe behandeln, alle cubischen Kreistheilungskörper, also alle Kreistheilungsperioden, die einer rationalen Gleichung $3^{\text{ten}}$ Grades genügen, aufzufinden, und zwar so, dass von mehreren Ausdrücken, deren dieselbe Grösse fähig ist, der einfachste gesetzt wird.

Wir haben hier nur eine Invariante $i_1 = 3$, und nach den Vorschriften in der Zusammenstellung des \S.~21 dürfen in $m$ aufgenommen werden der Factor $9$, nicht aber $3$, ferner Primzahlen in beliebiger Menge von der Form $6N+1$, also die Primzahlen $7, 13, 19, 31, 37, 43, 61, 67, 73, 79, 97\ldots$, und so bekommen wir also als geeignete Werthe von $m$ unter $200$:
\[
\begin{aligned}
m &= 7,\ 9,\ 13,\ 19,\ 31,\ 37,\ 43,\ 61,\ 63,\ 67,\ 73,\ 79,\ 91,\ 97,\ 103,\ 109, \\
&\quad 117,\ 127,\ 133,\ 139,\ 151,\ 157,\ 163,\ 171,\ 181,\ 193,\ 199,
\end{aligned}
\]
unter denen die Zahlen mit mehreren Primtheilern durch fetten Druck ausgezeichnet sind. Die kleinste Zahl $m$, in der mehr als zwei Primzahlen aufgehen, ist $7 \cdot 9 \cdot 13 = 819$.

Verstehen wir also unter $q_1, q_2, \ldots$ die in $m$ aufgehenden Primzahlen, so sind die Indexmoduln $c_1, c_2, \ldots$, wenn $m$ ungerade ist:
\[
c_1 = q_1 - 1, \quad c_2 = q_2 - 1, \quad \ldots
\]
oder wenn $m$ durch $9$ theilbar ist:
\[
c_1 = 6, \quad c_2 = q_2 - 1, \quad \ldots
\]
Es ist also, wenn $q_1$ eine Primzahl von der Form $6N+1$ ist, $\varphi(q_1) = q_1-1$ durch $6$, also durch $3$ theilbar, und wenn $m$ durch $9$ theilbar ist, $\varphi(9) = 6$ durch $3$ theilbar.

Es ist also in allen Fällen $c_h$ durch $3$ theilbar, und die einzige Invariante $i_1 = 3$.

Wir haben also nach \S.~21, (11):
\[
\delta_{1,h} = 3, \quad i_{1,h} = 1, \quad c_{1,h} = \frac{c_h}{3},
\]
und die eine Zahl $\gamma_{1,h}$ nach \S.~21, (17):
\begin{equation}\tag{1}
\gamma_{1,h} \equiv \frac{c_h}{3} \pmod{c_h}.
\end{equation}

Die Gruppe $\mathfrak{B}$ hat die eine Basis $g_1$ mit dem Grade $3$, und wenn $\beta_1, \beta_2, \ldots \beta_u$ die Indices irgend eines Elementes $b$ von $\mathfrak{B}$ sind, so ist nach \S.~21, (8):
\begin{equation}\tag{2}
\beta_h \equiv \gamma_{1,h} z_1 \pmod{c_h},
\end{equation}
worin $z_1 = 0, 1, 2$ ist.

Die Gruppe $\mathfrak{A}$ besteht aus allen den Zahlen, deren Indices die Congruenz \S.~22, (3) befriedigen, die hier lautet:
\begin{equation}\tag{3}
\frac{c_1}{3} \alpha_1 + \frac{c_2}{3} \alpha_2 + \cdots + \frac{c_u}{3} \alpha_u \equiv 0 \pmod{3}.
\end{equation}

Um die Gleichung für die Periode $\eta$ zu bilden, setzen wir
\begin{equation}\tag{4}
\eta = \sum_{a} r^a,
\end{equation}
worin $r$ eine primitive $m^{\text{te}}$ Einheitswurzel ist, und $a$ die Zahlen der Gruppe $\mathfrak{A}$ durchläuft. Die conjugirten Werthe von $\eta$ erhalten wir, wenn wir in (4) $r$ durch $r^b$ ersetzen, worin $b$ die Zahlen der Gruppe $\mathfrak{B}$ durchläuft.

Die Discriminante der cubischen Gleichung für $\eta$ ist nach den früheren Formeln
\[
D = m^2 \cdot B^2,
\]
worin $B$ eine ganze Zahl ist.

\medskip\noindent\textbf{1.} Um die cubische Gleichung für $\eta$ zu bilden, bezeichnen wir mit $\omega$ eine primitive dritte Einheitswurzel und setzen
\begin{equation}\tag{5}
\xi = (\eta_0 + \omega \eta_1 + \omega^2 \eta_2)^3,
\end{equation}
worin $\eta_0, \eta_1, \eta_2$ die drei conjugirten Werthe von $\eta$ sind.

Dann ist $\xi$ eine Zahl des Körpers $R(\omega)$, also von der Form $a + b\omega$, worin $a, b$ rationale Zahlen sind.

Setzen wir zur Abkürzung
\begin{equation}\tag{6}
\theta = \eta_0 + \omega \eta_1 + \omega^2 \eta_2,
\end{equation}
so ist
\begin{equation}\tag{7}
\theta' = \eta_0 + \omega^2 \eta_1 + \omega \eta_2,
\end{equation}
und es ist
\begin{equation}\tag{8}
\xi = \theta^3, \quad \xi' = \theta'^3.
\end{equation}

Nun ist $\theta \theta' = \eta_0^2 + \eta_1^2 + \eta_2^2 - \eta_0\eta_1 - \eta_1\eta_2 - \eta_2\eta_0 = -3B\sqrt{m}$ oder eine ähnliche Form, je nach der Zahl $m$.

Wir wollen nun für die einfachsten Fälle die Rechnung durchführen. Die Gleichung für $\eta$ hat die Form
\begin{equation}\tag{9}
\eta^3 - 3m\eta - mA = 0,
\end{equation}
worin $A$ eine ganze Zahl ist.

Ist $m = q$ eine Primzahl von der Form $6N+1$, so erhalten wir nach den Formeln des ersten Bandes, \S.~172, wenn $\varepsilon$ eine primitive dritte Einheitswurzel ist:
\begin{equation}\tag{10}
a + b\varepsilon = (\varrho, \varepsilon),
\end{equation}
worin $(\varrho, \varepsilon)$ eine Gauss'sche Summe ist, und
\begin{equation}\tag{11}
4q = A^2 + 27B^2,
\end{equation}
wobei $A \equiv 1 \pmod{3}$ ist.

Die cubische Gleichung für $\eta$ wird dann
\begin{equation}\tag{12}
\eta^3 - 3q\eta - qA = 0,
\end{equation}
und ihre Discriminante ist
\begin{equation}\tag{13}
D = q^2 \cdot 27B^2 = 27q^2B^2.
\end{equation}

Ist $m = q_1 q_2$ ein Product von zwei verschiedenen Primzahlen von der Form $6N+1$, so hat man die Gleichung
\begin{equation}\tag{14}
\eta^3 - 3m\eta - mA = 0,
\end{equation}
worin $A$ durch die Formel
\begin{equation}\tag{15}
4m = A^2 + 27B^2
\end{equation}
bestimmt ist.

\medskip\noindent\textbf{2.} Um die Gruppen $\mathfrak{A}$ zu finden, hat man zwei Zahlen $x, y$ so zu bestimmen, dass $7x + 9y = 1$ wird, und dann
\[
r = 7^x \cdot 9^y
\]
zu setzen. Nimmt man z.B. $x = 4, y = -3$ an, so erhält man für $m = 63$ die beiden Werthe von $\eta$:
\[
\eta = (r + r^8)(r^9 + r^{-9}) + (r^7 + r^{-7})(r^{18} + r^{-18})
\]
oder wenn man die Perioden der $7^{\text{ten}}$ und $9^{\text{ten}}$ Einheitswurzeln einführt:
\[
\eta = \eta_1 \vartheta_1 + \eta_2 \vartheta_2,
\]
worin $\eta_1, \eta_2$ die Perioden der $7^{\text{ten}}$ und $\vartheta_1, \vartheta_2$ die Perioden der $9^{\text{ten}}$ Einheitswurzeln sind.

Die Gruppen können durch die Potenzen von $2$ und von $11$ dargestellt werden, und die beiden Gruppen sind
\[
\mathfrak{A} = \{1, 8, 11, 25, 5, 23\} \quad \text{und} \quad \mathfrak{A}' = \{1, 8, 2, 16, 4, 32\}.
\]

\medskip\noindent\textbf{3.} Die cubischen Gleichungen, deren Bildungsgesetze wir jetzt kennen gelernt haben, enthalten mit ihren Tschirnhausen-Transformationen alle cubischen Kreistheilungsgleichungen. Aber diese Gleichungen sind auch alle von einander verschieden und können nicht durch Tschirnhausen-Transformation in einander übergeführt werden, weil sie zu primären Theilern von verschiedenen vollen Kreistheilungskörpern gehören, oder, wenn sie in demselben vollen Kreistheilungskörper enthalten sind, verschiedene Gruppen haben.

Man kann freilich noch aus anderen Einheitswurzeln Kreistheilungsperioden bilden, die zu cubischen Gleichungen führen. Diese Perioden können aber durch niedrigere Einheitswurzeln ausgedrückt werden, und sind nicht primär. So erhält man z.B., wenn $r$ eine $35^{\text{te}}$ Einheitswurzel ist, eine Periode von $8$ Gliedern, deren Exponenten aus den Potenzen von $-8$ gebildet sind:
\[
\eta = r + r^{-8} + r^{8} + r^{-1} + r^{4} + r^{-4} + r^{2} + r^{-2},
\]
die also auch einer cubischen Gleichung genügt. Es ist aber
\[
r + r^{7} + r^{-7} + r^{4} + r^{-4} + r^{2} + r^{-2} = 0,
\]
und wenn man diese Gleichung mit $(r^{15} + r^{-15})$ multiplicirt:
\[
r^{15} + r^{16} + r^{13} + r^{18} + r^{8} + r^{22} + r^{13} + r^{-3} = 0,
\]
also
\[
\eta = -(r^{15} + r^{-15}),
\]
was unter den Perioden der $7^{\text{ten}}$ Einheitswurzeln schon vorkommt. Diese Erscheinung erklärt sich daraus, dass $m = 35$ nicht die in \S.~21 verlangte Eigenschaft hat, dass $q-1$ für alle in $m$ aufgehenden Primzahlen $q$ durch $3$ theilbar ist.

\end{document}

